
\section{Methodology}


\subsection{Literature Research}

Initial scoping was guided by domain expert Marvyn Gulina \cite{Gulina23Koopman, Gulina17Chaos} and foundational literature \cite{Jaeger10ESN}. To establish theoretical framework, the seminal papers defining the Echo State Network \cite{Jaeger10ESN} paradigm and the work of University of Bordeaux \cite{Trouvain2020ReservoirPy} were analyzed first. 

\vspace{0.5cm}

The literature review was conducted formally using primarily Google Scholar. To ensure a reproducible search process,  my specific query was "reservoir computing real-time audio". 

\vspace{0.5cm}

Following this, the research focused on peer-reviewed literature addressing the problem of real-time distortion correction. This targeted search highlighted \textbf{RC} applications, notably the use of cascaded DLRC for guitar amplifier emulation \cite{Keuninckx2017RealTime}. By analyzing the reference lists of these highly relevant papers, other major architectural update within the RC paradigm were identified, such as the single dynamical node \cite{Appeltant2011Information} and deep reservoir model \cite{Gallicchio2017}. Finally, subsequent searches were conducted to find recent models capable of handling complex input streams with reduced computational overhead, such as deep time-delay models \cite{goldmann2020deep}.

\subsection{Literature Selection}

The selection of the literature was driven by several factors. First, I prioritized papers directly relevant to my field and research question, such as \cite{Keuninckx2017RealTime}. After this, I focused on foundational papers that introduced major structural modifications to the Reservoir Computing paradigm. The primary search engine I used here is \textbf{Google Scholar}, targeting reputable sources like \textit{Nature Communications}, \textit{Springer Nature}, \textit{ScienceDirect}, \textit{IEEE}. I also considered paper with at least 20 citation counts or those frequently referenced within the articles I analyzed.

\vspace{0.5cm}

Certain seminal articles were obvious choices for this literature review, such as \cite{Appeltant2011Information}, which had accumulated 2,133 citations, and \cite{maass2002realtime}, with 5,516 citations at the time of research.

\vspace{0.5cm}

After this I search for more recent papers under the \textbf{RC} paradigm with a focus on newer models that can compute real-time non linear distortion on more complex input stream or in faster rate such as \cite{goldmann2020deep, Gauthier2021, peixeiro2022time}

\vspace{0.5cm}

My primary exclusion criterion is paper covering hardware and electronics implementation of \textbf{RC}. I encountered numerous articles discussing laser dynamics, satellite distortion, and electronic oscillators. Lacking a formal background in electronics, I chose to focus exclusively on software-based approaches to remain within my area of expertise.

\vspace{0.5cm}

My second exclusion criterion was based on publication dates, as I aimed to prioritize modern models within the RC paradigm. Consequently, I excluded papers older than the year 2000, which marks the beginning of the Reservoir Computing era with the work of Jaeger \cite{Jaeger10ESN} and Maass \cite{maass2002realtime}.

\vspace{0.5cm}


The third exclusion criterion applied to biological studies. While Reservoir Computing falls under Artificial Intelligence, many papers discuss parallels with biological neurons. For obvious reasons, a biological review is beyond the scope of this work. 

\vspace{0.5cm}

Finally, as Reservoir Computing involves complex mathematical frameworks, the equations presented in this paper are sourced directly from the referenced literature to ensure technical accuracy and rigor.

\subsection{Extraction and analysis}

The primary extraction and analysis methodology used in this review began with a reading of the articles. This initial step was followed by a study of the mathematical foundations underlying each computational paradigm. Then a comparative analysis of their experimental results.

\subsection{Use of Generative AI}

In accordance with institutional and academic integrity guidelines, this section details the use of AI. The Gemini model (Google) was strictly used to improve the quality of the text. This included orthographic and grammatical corrections, as well as stylistic reprocessing of certain passages to achieve a more formal and academic tone. The AI was used exclusively as a writing aid. All original ideas, research content, and analysis were developed solely by the author, and no content was generated directly from it.